% ================= Feature 002 (Lê Sĩ Nhật Quang) =================
\section{Feature 002: Add a new course}

\subsection{Mô tả và hướng dẫn sử dụng}
Quản trị viên tạo khóa học mới tại \textbf{Site administration $>$ Courses $>$ Add a new course}
(URL: \texttt{<BASE\_URL>/course/edit.php?category=0}), hoặc từ \textbf{Manage courses and categories $>$ Create new course}.
Form gồm các trường Course full name, Course short name, Course category, Course visibility, Course start date,
Course end date (có checkbox \textit{Enable}) và Course ID number; ngoài ra còn các nhóm thiết lập mở rộng
(Description, Course format, Appearance,\dots) giữ giá trị mặc định.

\textbf{Các bước sử dụng:}
\begin{enumerate}[noitemsep]
  \item Đăng nhập bằng tài khoản admin.
  \item Mở trang Add a new course.
  \item Nhập Course full name, Course short name; chọn Course category bằng ô tìm kiếm (autocomplete).
  \item Chọn Course visibility (Show/Hide), đặt Course start date; bật/tắt và đặt Course end date; nhập Course ID number (tùy chọn).
  \item Bấm \textbf{Save and display} (mở trang khóa học), \textbf{Save and return} (về trang quản lý) hoặc \textbf{Cancel}.
  \item Nếu dữ liệu không hợp lệ, form hiển thị lỗi tại trường tương ứng và không tạo khóa học.
\end{enumerate}

\subsection{Use case và activity diagram}
\begin{table}[H]
\centering\small
\begin{tabular}{|L{3.6cm}|L{11.4cm}|}
\hline
Tên & UC-002: Tạo khóa học mới \\ \hline
Tác nhân & Quản trị viên (Admin) \\ \hline
Tiền điều kiện & Admin đã đăng nhập; tồn tại ít nhất một category (``Category 1''). \\ \hline
Hậu điều kiện (thành công) & Khóa học mới được lưu trong category đã chọn và hiển thị trong trang quản lý khóa học. \\ \hline
Luồng chính & Mở form $\rightarrow$ nhập dữ liệu hợp lệ $\rightarrow$ Save and display $\rightarrow$ hệ thống tạo khóa học và mở trang khóa học. \\ \hline
Luồng thay thế & A1: bấm Save and return $\rightarrow$ tạo khóa học, quay về trang quản lý. A2: bấm Cancel $\rightarrow$ không tạo khóa học. \\ \hline
Luồng ngoại lệ & E1: thiếu trường bắt buộc (full name, short name, category). E2: short name đã tồn tại. E3: ID number đã tồn tại.
E4: end date không lớn hơn start date. E5: vượt độ dài tối đa. Mọi ngoại lệ: báo lỗi, giữ form để sửa và gửi lại. \\ \hline
\end{tabular}
\caption{Đặc tả use case UC-002}
\end{table}

\begin{figure}[H]
\centering
\begin{tikzpicture}[>=Stealth,
  act/.style={rectangle, rounded corners=4pt, draw, minimum width=3.6cm, minimum height=0.8cm, align=center, font=\small, fill=blue!5},
  dec/.style={diamond, aspect=2.2, draw, align=center, font=\scriptsize, inner sep=1pt, fill=yellow!15},
  st/.style={circle, fill=black, minimum size=0.35cm},
  en/.style={circle, draw, double, minimum size=0.35cm}]
\node[st] (s) {};
\node[act, below=0.5cm of s] (a1) {Mở trang Add a new course};
\node[act, below=0.5cm of a1] (a2) {Nhập/sửa dữ liệu form};
\node[dec, below=0.6cm of a2] (d0) {Bấm\\Cancel?};
\node[act, below=0.7cm of d0] (a3) {Bấm Save and display\\hoặc Save and return};
\node[dec, below=0.6cm of a3] (d1) {Đủ trường\\bắt buộc, đúng\\độ dài?};
\node[dec, below=0.7cm of d1] (d2) {Short name/\\ID number\\không trùng?};
\node[dec, below=0.7cm of d2] (d3) {End date tắt\\hoặc $>$ start?};
\node[dec, below=0.7cm of d3] (d4) {Save and\\display?};
\node[act, below left=0.7cm and -0.6cm of d4] (a4) {Mở trang khóa học};
\node[act, below right=0.7cm and -0.6cm of d4] (a5) {Về trang quản lý};
\node[en, below=2.2cm of d4] (e) {};
\node[act, right=2.3cm of d0] (c1) {Về trang quản lý,\\không tạo course};
\node[act, right=2.3cm of d1] (er1) {Lỗi Missing/\\Maximum of\dots};
\node[act, right=2.3cm of d2] (er2) {Lỗi ``already used''};
\node[act, right=2.3cm of d3] (er3) {Lỗi end date};
\draw[->] (s)--(a1); \draw[->] (a1)--(a2); \draw[->] (a2)--(d0);
\draw[->] (d0)--node[right,font=\scriptsize]{Không}(a3);
\draw[->] (d0)--node[above,font=\scriptsize]{Có}(c1);
\draw[->] (a3)--(d1);
\draw[->] (d1)--node[right,font=\scriptsize]{Có}(d2);
\draw[->] (d1)--node[above,font=\scriptsize]{Không}(er1);
\draw[->] (d2)--node[right,font=\scriptsize]{Có}(d3);
\draw[->] (d2)--node[above,font=\scriptsize]{Không}(er2);
\draw[->] (d3)--node[right,font=\scriptsize]{Có}(d4);
\draw[->] (d3)--node[above,font=\scriptsize]{Không}(er3);
\draw[->] (d4)-|node[above,font=\scriptsize,pos=0.25]{Có}(a4);
\draw[->] (d4)-|node[above,font=\scriptsize,pos=0.25]{Không}(a5);
\draw[->] (a4)|-(e); \draw[->] (a5)|-(e);
\coordinate (R) at ($(er1.east)+(0.5,0)$);
\coordinate (RC) at ($(R)+(0.7,0)$);
\draw (er1.east)--(er1.east -| R); \draw (er2.east)--(er2.east -| R); \draw (er3.east)--(er3.east -| R);
\draw[->] (er3.east -| R)--(a2.east -| R)--(a2.east);
\node[font=\scriptsize, fill=white, inner sep=1pt] at ($(a2.east -| R)!0.5!(a2.east)$) {sửa lỗi};
\draw[->] (c1.east)--(c1.east -| RC)|-(e.east);
\end{tikzpicture}
\caption{Activity diagram của UC-002 (các nhánh lỗi quay lại bước nhập dữ liệu; Cancel kết thúc không tạo khóa học)}
\end{figure}

\subsection{Các biến đầu vào}
\begin{table}[H]
\centering\small
\begin{tabular}{|L{2.8cm}|L{1.8cm}|L{5.0cm}|L{4.4cm}|}
\hline
\rowcolor{gray!15}\textbf{Biến} & \textbf{Bắt buộc?} & \textbf{Ràng buộc} & \textbf{Thông báo lỗi} \\ \hline
Course full name & Có & 1--254 ký tự; chuỗi chỉ có khoảng trắng coi như rỗng & Missing full name / Maximum of 254 characters \\ \hline
Course short name & Có & 1--100 ký tự; không trùng khóa học khác & Missing short name / Maximum of 100 characters / Short name is already used for another course (\dots) \\ \hline
Course category & Có & Chọn một category tồn tại & You must supply a value here. \\ \hline
Course visibility & Có (mặc định Show) & Show / Hide & Không có giá trị sai qua giao diện \\ \hline
Course start date & Có & Ngày giờ chọn từ dropdown (năm 1900--2050) & Không có giá trị sai qua giao diện \\ \hline
Course end date & Không (checkbox Enable) & Nếu Enable: phải $>$ start date & The course end date must be after the start date. \\ \hline
Course ID number & Không & 0--100 ký tự; nếu có thì không trùng & ID number is already used for another course (\dots) \\ \hline
\end{tabular}
\caption{Các biến đầu vào của Feature 002}
\end{table}
\textit{Ghi chú:} giới hạn độ dài trong bảng là giới hạn \emph{thiết kế} (đặc tả ban đầu: thuộc tính \texttt{maxlength} và quy tắc
kiểm tra phía trình duyệt của form). Khi chạy thực tế bằng Katalon Recorder (lệnh \texttt{type} gán giá trị trực tiếp cho ô nhập),
Moodle demo \textbf{chấp nhận} full name chỉ có khoảng trắng, full name 255 ký tự, short name 101 ký tự và ID number 101 ký tự;
xem mục~\ref{sec:f002-result}.

\subsection{Kỹ thuật 1: Phân lớp tương đương (ECP)}
\textbf{Cách áp dụng:} chia miền giá trị mỗi biến thành lớp hợp lệ (V) và không hợp lệ (I). Hai test case ``tất cả hợp lệ''
(TC-002-001, 002) phủ hết các lớp V (mỗi lựa chọn Show/Hide, end date bật/tắt, ID có/không); mỗi lớp I có một test case riêng,
các trường khác giữ hợp lệ.

\begin{table}[H]
\centering\footnotesize
\begin{tabular}{|l|L{5.0cm}|L{7.4cm}|}
\hline
\rowcolor{gray!15}\textbf{Biến} & \textbf{Lớp hợp lệ (TC)} & \textbf{Lớp không hợp lệ (TC)} \\ \hline
Full name & V1: 1--254 ký tự (001, 002) & I1: rỗng (003); I2: chỉ khoảng trắng (004); I3: 255 ký tự (005) \\ \hline
Short name & V2: 1--100 ký tự, chưa tồn tại (001, 002) & I4: rỗng (006); I5: 101 ký tự (007); I6: trùng EXIST01 (008) \\ \hline
Category & V3: chọn Category 1 (001, 002) & I7: No selection (009) \\ \hline
Visibility & V4: Show (001); V5: Hide (002) & -- \\ \hline
End date & V6: tắt Enable (002); V7: bật, end $>$ start (001) & I8: end = start (010); I9: end $<$ start (011) \\ \hline
ID number & V8: rỗng (002); V9: 1--100 ký tự, chưa tồn tại (001) & I10: 101 ký tự (012); I11: trùng ID-EXIST (013) \\ \hline
\end{tabular}
\caption{Lớp tương đương của Feature 002 (ID rút gọn: 003 nghĩa là TC-002-003)}
\end{table}
\textit{Sau khi chạy:} các lớp I2, I3, I5, I10 được thiết kế là \emph{không hợp lệ} nhưng hệ thống thực tế chấp nhận
(khóa học vẫn được tạo). Trong file \texttt{.xlsx}, mô tả của TC-002-004, 005, 007, 012 được cập nhật thành ``lớp hợp lệ''
và kết quả mong đợi được sửa theo hành vi thực tế; sự khác biệt này được ghi nhận là phát hiện kiểm thử.

\subsection{Kỹ thuật 2: Phân tích giá trị biên (BVA)}
\textbf{Cách áp dụng:} chọn min$-1$, min, min$+1$, max$-1$, max, max$+1$ cho mỗi biến có miền liên tục. Với ngày, ``độ dài''
là khoảng cách end -- start với bước nhỏ nhất là 1 phút (đơn vị nhỏ nhất của dropdown). Giá trị không hợp lệ trùng ECP được tham chiếu lại.

\begin{table}[H]
\centering\footnotesize
\begin{tabular}{|l|c|c|c|c|c|c|}
\hline
\rowcolor{gray!15}\textbf{Biến} & \textbf{min$-1$} & \textbf{min} & \textbf{min$+1$} & \textbf{max$-1$} & \textbf{max} & \textbf{max$+1$} \\ \hline
Full name (1--254) & 0 (003) & 1 (014) & 2 (015) & 253 (016) & 254 (017) & 255 (005) \\ \hline
Short name (1--100) & 0 (006) & 1 (018) & 2 (019) & 99 (020) & 100 (021) & 101 (007) \\ \hline
ID number (0--100) & -- & 0 (002) & 1 (022) \;/ 2 (023) & 99 (024) & 100 (025) & 101 (012) \\ \hline
End $-$ start & $-1$ phút (026) & 0 (010)$^{*}$ & $+1$ phút (027) & \multicolumn{3}{c|}{không có cận trên (năm tối đa 2050)} \\ \hline
\end{tabular}
\caption{Giá trị biên của Feature 002 ($^{*}$end = start là giá trị \emph{không hợp lệ}: biên hợp lệ nhỏ nhất là +1 phút;
max$+1$ của full name, short name, ID number thực tế được chấp nhận, xem mục~\ref{sec:f002-result})}
\end{table}

\subsection{Kỹ thuật 3: Bảng quyết định (Decision table)}
\textbf{Điều kiện:} C1 -- full name hợp lệ; C2 -- short name hợp lệ và không trùng; C3 -- đã chọn category;
C4 -- bật end date; C5 -- end date $>$ start date.
\textbf{Hành động:} A1 -- tạo khóa học; A2 -- lỗi full name; A3 -- lỗi short name; A4 -- lỗi category; A5 -- lỗi end date.
Các lỗi bắt buộc (A2--A4) được kiểm tra phía trình duyệt nên hiển thị đồng thời; lỗi trùng và lỗi ngày được kiểm tra phía server.

\begin{table}[H]
\centering\footnotesize
\begin{tabular}{|l|c|c|c|c|c|c|c|}
\hline
\rowcolor{gray!15}\textbf{Luật} & R1 & R2 & R3 & R4 & R5 & R6 & R7 \\ \hline
C1 Full name hợp lệ & Y & Y & Y & N & Y & Y & Y \\ \hline
C2 Short name hợp lệ & Y & Y & Y & N & N (trùng) & N (rỗng) & Y \\ \hline
C3 Có category & Y & Y & Y & N & Y & Y & N \\ \hline
C4 Bật end date & N & Y & Y & -- & Y & -- & -- \\ \hline
C5 End $>$ start & -- & Y & N & -- & N & -- & -- \\ \hline
A1 Tạo khóa học & X & X & & & & & \\ \hline
A2 Lỗi full name & & & & X & & & \\ \hline
A3 Lỗi short name & & & & X & X & X & \\ \hline
A4 Lỗi category & & & & X & & & X \\ \hline
A5 Lỗi end date & & & X & & X & & \\ \hline
\textbf{Test case} & 028 & 029 & 030 & 031 & 032 & 006 & 009 \\ \hline
\end{tabular}
\caption{Decision table của Feature 002 (R6, R7 tái sử dụng test case ECP)}
\end{table}

\subsection{Kỹ thuật 4: Use-case testing}
\begin{table}[H]
\centering\footnotesize
\begin{tabular}{|c|L{7.4cm}|L{3.2cm}|c|}
\hline
\rowcolor{gray!15}\textbf{Path} & \textbf{Đường đi trong activity diagram} & \textbf{Loại} & \textbf{Test case} \\ \hline
P1 & Mở form $\rightarrow$ nhập hợp lệ $\rightarrow$ Save and display $\rightarrow$ trang khóa học & Luồng chính & 033 \\ \hline
P2 & Mở form $\rightarrow$ nhập hợp lệ $\rightarrow$ Save and return $\rightarrow$ trang quản lý & Luồng thay thế A1 & 034 \\ \hline
P3 & Nhập dữ liệu $\rightarrow$ Cancel $\rightarrow$ không tạo khóa học & Luồng thay thế A2 & 035 \\ \hline
P4 & Gửi form rỗng $\rightarrow$ lỗi bắt buộc $\rightarrow$ nhập lại $\rightarrow$ thành công & Ngoại lệ E1 & 036 \\ \hline
P5 & Short name trùng $\rightarrow$ lỗi ``already used'' $\rightarrow$ đổi $\rightarrow$ thành công & Ngoại lệ E2 & 037 \\ \hline
P6 & End date $<$ start $\rightarrow$ lỗi end date $\rightarrow$ sửa $\rightarrow$ thành công & Ngoại lệ E4 & 038 \\ \hline
\end{tabular}
\caption{Scenario và test case use-case testing của Feature 002 (E3 đã phủ bởi TC-002-013; E5 được kiểm tra bởi TC-002-005, 007, 012 nhưng thực tế hệ thống không báo lỗi, xem mục~\ref{sec:f002-result})}
\end{table}

\subsection{Kết quả thực thi và phát hiện}\label{sec:f002-result}
Toàn bộ 38 test case đã được ghi và chạy bằng Katalon Recorder trên Moodle demo (file \texttt{F002-Add a new course.krecorder}).
Kết quả mong đợi trong file \texttt{.xlsx} được cập nhật theo hành vi thực tế quan sát được. Các test case có kết quả khác đặc tả ban đầu:

\begin{table}[H]
\centering\footnotesize
\begin{tabular}{|l|L{3.3cm}|L{4.3cm}|L{4.6cm}|}
\hline
\rowcolor{gray!15}\textbf{Test case} & \textbf{Đầu vào} & \textbf{Kết quả dự kiến ban đầu} & \textbf{Kết quả thực tế} \\ \hline
TC-002-004 & Full name = 3 khoảng trắng & Lỗi ``Missing full name'' & Khóa học được tạo với full name là khoảng trắng \\ \hline
TC-002-005 & Full name 255 ký tự & Lỗi ``Maximum of 254 characters'' hoặc bị cắt còn 254 & Khóa học được tạo với full name \texttt{Faaa\dots} \\ \hline
TC-002-007 & Short name 101 ký tự & Lỗi ``Maximum of 100 characters'' & Khóa học được tạo \\ \hline
TC-002-012 & ID number 101 ký tự & Lỗi ``Maximum of 100 characters'' hoặc bị cắt còn 100 & Khóa học được tạo \\ \hline
\end{tabular}
\caption{Khác biệt giữa đặc tả ban đầu và hành vi thực tế của Feature 002}
\end{table}

\textbf{Nhận xét:}
\begin{itemize}[noitemsep]
  \item Giới hạn độ dài của form chỉ được áp dụng ở phía trình duyệt (\texttt{maxlength}, kiểm tra JavaScript). Lệnh \texttt{type}
        của Katalon gán giá trị trực tiếp nên vượt qua được, và phía server không chặn vì cột dữ liệu trong CSDL dài hơn
        giới hạn hiển thị trên form.
  \item Full name chỉ có khoảng trắng vẫn được lưu: khóa học có tên ``trống'' gây khó nhận diện trong danh sách. Đây là điểm
        nên báo cáo như một lỗi tiềm ẩn (mức độ thấp).
  \item Các test case còn lại cho kết quả đúng như mong đợi: các lỗi bắt buộc (Missing full name, Missing short name, You must supply a value here.),
        lỗi trùng short name/ID number và lỗi end date đều xuất hiện đúng trường.
\end{itemize}

\subsection{Danh sách test case của Feature 002}
Tổng cộng \textbf{38 test case}: ECP 001--013, BVA 014--027, Decision table 028--032, Use-case 033--038.
Kết quả mong đợi chung: test case hợp lệ -- khóa học được tạo và hiển thị trong \textit{My courses} với đúng full name;
test case không hợp lệ -- vẫn ở form Add a new course, hiển thị lỗi tại trường tương ứng và không có khóa học mới.
Riêng TC-002-004, 005, 007, 012 có kết quả mong đợi theo hành vi thực tế (mục~\ref{sec:f002-result}).
Dữ liệu dùng tiền tố \texttt{TC002} để dễ dọn dẹp.

{\footnotesize
\begin{longtable}{|l|L{4.4cm}|c|L{6.4cm}|}
\hline
\rowcolor{gray!15}\textbf{ID} & \textbf{Tên test case} & \textbf{Kỹ thuật} & \textbf{Mô tả} \\ \hline
\endhead
TC-002-001 & Tạo course với dữ liệu hợp lệ (Show, có end date, có ID) & ECP & Tất cả các trường thuộc lớp hợp lệ V1, V2, V3, V4, V7, V9. \\ \hline
TC-002-002 & Tạo course Hide, tắt end date, không ID & ECP & Lớp hợp lệ V1, V2, V3, V5, V6, V8. \\ \hline
TC-002-003 & Full name rỗng & ECP & Lớp không hợp lệ I1: để trống Course full name. \\ \hline
TC-002-004 & Full name chỉ có khoảng trắng & ECP & Lớp hợp lệ I2: Course full name gồm 3 khoảng trắng. \\ \hline
TC-002-005 & Full name 255 ký tự & ECP/BVA & Lớp hợp lệ I3: Course full name dài 255 ký tự (tối đa 254). \\ \hline
TC-002-006 & Short name rỗng & ECP & Lớp không hợp lệ I4: để trống Course short name. \\ \hline
TC-002-007 & Short name 101 ký tự & ECP/BVA & Lớp hợp lệ I5: Course short name dài 101 ký tự. \\ \hline
TC-002-008 & Short name đã tồn tại & ECP & Lớp không hợp lệ I6: Course short name trùng "EXIST01" của course có sẵn. \\ \hline
TC-002-009 & Không chọn Course category & ECP & Lớp không hợp lệ I7: Course category để "No selection". \\ \hline
TC-002-010 & End date bằng start date & ECP/BVA & Lớp không hợp lệ I8: end date = start date (1 January 2026 08:00). \\ \hline
TC-002-011 & End date trước start date & ECP & Lớp không hợp lệ I9: end date (1 January 2025 08:00) \textless{} start date. \\ \hline
TC-002-012 & ID number 101 ký tự & ECP/BVA & Lớp hợp lệ I10: Course ID number dài 101 ký tự (tối đa 100). \\ \hline
TC-002-013 & ID number đã tồn tại & ECP & Lớp không hợp lệ I11: Course ID number trùng "ID-EXIST". \\ \hline
TC-002-014 & Full name 1 ký tự & BVA & Biên hợp lệ của Course full name: độ dài 1. \\ \hline
TC-002-015 & Full name 2 ký tự & BVA & Biên hợp lệ của Course full name: độ dài 2. \\ \hline
TC-002-016 & Full name 253 ký tự & BVA & Biên hợp lệ của Course full name: độ dài 253. \\ \hline
TC-002-017 & Full name 254 ký tự & BVA & Biên hợp lệ của Course full name: độ dài 254. \\ \hline
TC-002-018 & Short name 1 ký tự & BVA & Biên hợp lệ của Course short name: độ dài 1. \\ \hline
TC-002-019 & Short name 2 ký tự & BVA & Biên hợp lệ của Course short name: độ dài 2. \\ \hline
TC-002-020 & Short name 99 ký tự & BVA & Biên hợp lệ của Course short name: độ dài 99. \\ \hline
TC-002-021 & Short name 100 ký tự & BVA & Biên hợp lệ của Course short name: độ dài 100. \\ \hline
TC-002-022 & ID number 1 ký tự & BVA & Biên hợp lệ của Course ID number: độ dài 1. \\ \hline
TC-002-023 & ID number 2 ký tự & BVA & Biên hợp lệ của Course ID number: độ dài 2. \\ \hline
TC-002-024 & ID number 99 ký tự & BVA & Biên hợp lệ của Course ID number: độ dài 99. \\ \hline
TC-002-025 & ID number 100 ký tự & BVA & Biên hợp lệ của Course ID number: độ dài 100. \\ \hline
TC-002-026 & End date = start date - 1 phút & BVA & Biên dưới không hợp lệ: end 1 January 2026 07:59, start 1 January 2026 08:00. \\ \hline
TC-002-027 & End date = start date + 1 phút & BVA & Biên hợp lệ nhỏ nhất: end 1 January 2026 08:01, start 1 January 2026 08:00. \\ \hline
TC-002-028 & DT-R1: Hợp lệ, end date tắt & DT & R1: C1=T, C2=T, C3=T, C4=F (C5 don't care) -\textgreater{} tạo course thành công. \\ \hline
TC-002-029 & DT-R2: Hợp lệ, end date \textgreater{} start & DT & R2: C1=T, C2=T, C3=T, C4=T, C5=T -\textgreater{} tạo course thành công. \\ \hline
TC-002-030 & DT-R3: End date không hợp lệ & DT & R3: C1=T, C2=T, C3=T, C4=T, C5=F -\textgreater{} lỗi end date. \\ \hline
TC-002-031 & DT-R4: Thiếu toàn bộ trường bắt buộc & DT & R4: C1=F, C2=F, C3=F -\textgreater{} hiển thị đồng thời 3 lỗi. \\ \hline
TC-002-032 & DT-R5: Short name trùng và end date sai & DT & R5: C1=T, C2=F (trùng), C3=T, C4=T, C5=F -\textgreater{} hiển thị cả 2 lỗi server. \\ \hline
TC-002-033 & UC-P1: Luồng chính - Save and display & UC & Basic flow: nhập dữ liệu hợp lệ, Save and display. \\ \hline
TC-002-034 & UC-P2: Luồng thay thế - Save and return & UC & Alternative flow A1: lưu và quay về trang quản lý. \\ \hline
TC-002-035 & UC-P3: Hủy tạo course & UC & Alternative flow A2: nhập dữ liệu rồi Cancel. \\ \hline
TC-002-036 & UC-P4: Bỏ trống rồi nhập lại & UC & Exception E1 -\textgreater{} sửa -\textgreater{} thành công. \\ \hline
TC-002-037 & UC-P5: Short name trùng rồi đổi & UC & Exception E2 -\textgreater{} sửa -\textgreater{} thành công. \\ \hline
TC-002-038 & UC-P6: End date sai rồi sửa & UC & Exception E4 -\textgreater{} sửa -\textgreater{} thành công. \\ \hline
\end{longtable}
}


\subsection{Test case mẫu của Feature 002}
Sáu test case dưới đây được trích nguyên từ file \texttt{.xlsx} (sheet \textit{Add a new course}); các test case còn lại cùng khuôn.

\bigskip\noindent\textbf{TC-002-001 -- Tạo course với dữ liệu hợp lệ (Show, có end date, có ID)}\\
\noindent\textit{Precondition:} Moodle demo (Moodle 5.1, Docker) đang chạy tại http:/\allowbreak{}/\allowbreak{}127.0.0.1:9501 giao diện Moodle và trình duyệt là English; đã đăng nhập bằng admin (admin/\allowbreak{}admin). Đã có category "Category 1". Script Setup đã tạo sẵn course Full name "Existing Course", Short name "EXIST01", ID number "ID-EXIST" trong "Category 1". Chưa tồn tại course nào có Short name/\allowbreak{}ID number dùng trong test case này (script Cleanup xóa các course TC002-* sau khi chạy).
\smallskip

{\footnotesize
\begin{longtable}{|c|L{3.0cm}|L{5.1cm}|L{5.6cm}|}
\hline
\rowcolor{gray!15}\textbf{\#} & \textbf{Bước} & \textbf{Mô tả} & \textbf{Kết quả mong đợi} \\ \hline
\endhead
1 & Mở trang Add a new course & Vào Site administration \textgreater{} Courses \textgreater{} Add a new course (hoặc mở \textless{}BASE\_\allowbreak{}URL\textgreater{}/\allowbreak{}course/\allowbreak{}edit.php?category=0). & Form "Add a new course" hiển thị: Course full name, Course short name, Course category (No selection), Course visibility, Course start date, Course end date, Course ID number. \\ \hline
2 & Nhập Course full name & Nhập "TC002 Course 001" vào ô Course full name. & Ô Course full name hiển thị giá trị đã nhập. \\ \hline
3 & Nhập Course short name & Nhập "TC002-001" vào ô Course short name. & Ô Course short name hiển thị giá trị đã nhập. \\ \hline
4 & Chọn Course category & Click ô Search của Course category, gõ "Category 1" và chọn "Category 1" trong danh sách. & Course category hiển thị "Category 1". \\ \hline
5 & Chọn Course visibility & Chọn "Show" trong dropdown Course visibility. & Dropdown hiển thị "Show". \\ \hline
6 & Đặt Course start date & Chọn ngày/\allowbreak{}tháng/\allowbreak{}năm/\allowbreak{}giờ/\allowbreak{}phút = 1 January 2026 08:00. & Course start date hiển thị 1 January 2026 08:00. \\ \hline
7 & Đặt Course end date & Tick "Enable" của Course end date, chọn 1 January 2027 08:00. & Course end date hiển thị 1 January 2027 08:00. \\ \hline
8 & Nhập Course ID number & Nhập "ID-002-001" vào ô Course ID number. & Ô Course ID number hiển thị giá trị đã nhập. \\ \hline
9 & Lưu course & Click nút "Save and display". & Không có thông báo lỗi; course được tạo và trang course hiển thị tiêu đề "TC002 Course 001"; breadcrumb chứa short name "TC002-001". \\ \hline
\end{longtable}
}

\bigskip\noindent\textbf{TC-002-004 -- Full name chỉ có khoảng trắng}\\
\noindent\textit{Precondition:} Moodle demo (Moodle 5.1, Docker) đang chạy tại http:/\allowbreak{}/\allowbreak{}127.0.0.1:9501 giao diện Moodle và trình duyệt là English; đã đăng nhập bằng admin (admin/\allowbreak{}admin). Đã có category "Category 1". Script Setup đã tạo sẵn course Full name "Existing Course", Short name "EXIST01", ID number "ID-EXIST" trong "Category 1". Chưa tồn tại course nào có Short name/\allowbreak{}ID number dùng trong test case này (script Cleanup xóa các course TC002-* sau khi chạy).
\smallskip

{\footnotesize
\begin{longtable}{|c|L{3.0cm}|L{5.1cm}|L{5.6cm}|}
\hline
\rowcolor{gray!15}\textbf{\#} & \textbf{Bước} & \textbf{Mô tả} & \textbf{Kết quả mong đợi} \\ \hline
\endhead
1 & Mở trang Add a new course & Vào Site administration \textgreater{} Courses \textgreater{} Add a new course (hoặc mở \textless{}BASE\_\allowbreak{}URL\textgreater{}/\allowbreak{}course/\allowbreak{}edit.php?category=0). & Form "Add a new course" hiển thị: Course full name, Course short name, Course category (No selection), Course visibility, Course start date, Course end date, Course ID number. \\ \hline
2 & Nhập Course full name & Nhập 3 ký tự khoảng trắng vào ô Course full name. & Ô Course full name hiển thị giá trị đã nhập. \\ \hline
3 & Nhập Course short name & Nhập "TC002-004" vào ô Course short name. & Ô Course short name hiển thị giá trị đã nhập. \\ \hline
4 & Chọn Course category & Click ô Search của Course category, gõ "Category 1" và chọn "Category 1" trong danh sách. & Course category hiển thị "Category 1". \\ \hline
5 & Chọn Course visibility & Chọn "Show" trong dropdown Course visibility. & Dropdown hiển thị "Show". \\ \hline
6 & Đặt Course start date & Chọn ngày/\allowbreak{}tháng/\allowbreak{}năm/\allowbreak{}giờ/\allowbreak{}phút = 1 January 2026 08:00. & Course start date hiển thị 1 January 2026 08:00. \\ \hline
7 & Đặt Course end date & Tick "Enable" của Course end date, chọn 1 January 2027 08:00. & Course end date hiển thị 1 January 2027 08:00. \\ \hline
8 & Nhập Course ID number & Để trống (xóa hết nội dung) vào ô Course ID number. & Ô Course ID number hiển thị giá trị đã nhập. \\ \hline
9 & Lưu course & Click nút "Save and display". & Form được lưu với tên là 3 khoảng trắng \\ \hline
\end{longtable}
}

\bigskip\noindent\textbf{TC-002-008 -- Short name đã tồn tại}\\
\noindent\textit{Precondition:} Moodle demo (Moodle 5.1, Docker) đang chạy tại http:/\allowbreak{}/\allowbreak{}127.0.0.1:9501 giao diện Moodle và trình duyệt là English; đã đăng nhập bằng admin (admin/\allowbreak{}admin). Đã có category "Category 1". Script Setup đã tạo sẵn course Full name "Existing Course", Short name "EXIST01", ID number "ID-EXIST" trong "Category 1". Chưa tồn tại course nào có Short name/\allowbreak{}ID number dùng trong test case này (script Cleanup xóa các course TC002-* sau khi chạy).
\smallskip

{\footnotesize
\begin{longtable}{|c|L{3.0cm}|L{5.1cm}|L{5.6cm}|}
\hline
\rowcolor{gray!15}\textbf{\#} & \textbf{Bước} & \textbf{Mô tả} & \textbf{Kết quả mong đợi} \\ \hline
\endhead
1 & Mở trang Add a new course & Vào Site administration \textgreater{} Courses \textgreater{} Add a new course (hoặc mở \textless{}BASE\_\allowbreak{}URL\textgreater{}/\allowbreak{}course/\allowbreak{}edit.php?category=0). & Form "Add a new course" hiển thị: Course full name, Course short name, Course category (No selection), Course visibility, Course start date, Course end date, Course ID number. \\ \hline
2 & Nhập Course full name & Nhập "TC002 Course 008" vào ô Course full name. & Ô Course full name hiển thị giá trị đã nhập. \\ \hline
3 & Nhập Course short name & Nhập "EXIST01" vào ô Course short name. & Ô Course short name hiển thị giá trị đã nhập. \\ \hline
4 & Chọn Course category & Click ô Search của Course category, gõ "Category 1" và chọn "Category 1" trong danh sách. & Course category hiển thị "Category 1". \\ \hline
5 & Chọn Course visibility & Chọn "Show" trong dropdown Course visibility. & Dropdown hiển thị "Show". \\ \hline
6 & Đặt Course start date & Chọn ngày/\allowbreak{}tháng/\allowbreak{}năm/\allowbreak{}giờ/\allowbreak{}phút = 1 January 2026 08:00. & Course start date hiển thị 1 January 2026 08:00. \\ \hline
7 & Đặt Course end date & Tick "Enable" của Course end date, chọn 1 January 2027 08:00. & Course end date hiển thị 1 January 2027 08:00. \\ \hline
8 & Nhập Course ID number & Để trống (xóa hết nội dung) vào ô Course ID number. & Ô Course ID number hiển thị giá trị đã nhập. \\ \hline
9 & Lưu course & Click nút "Save and display". & Form không được lưu, vẫn ở trang Add a new course; hiển thị lỗi "Short name is already used for another course (Existing Course)". Không có course mới được tạo. \\ \hline
\end{longtable}
}

\bigskip\noindent\textbf{TC-002-010 -- End date bằng start date}\\
\noindent\textit{Precondition:} Moodle demo (Moodle 5.1, Docker) đang chạy tại http:/\allowbreak{}/\allowbreak{}127.0.0.1:9501 giao diện Moodle và trình duyệt là English; đã đăng nhập bằng admin (admin/\allowbreak{}admin). Đã có category "Category 1". Script Setup đã tạo sẵn course Full name "Existing Course", Short name "EXIST01", ID number "ID-EXIST" trong "Category 1". Chưa tồn tại course nào có Short name/\allowbreak{}ID number dùng trong test case này (script Cleanup xóa các course TC002-* sau khi chạy).
\smallskip

{\footnotesize
\begin{longtable}{|c|L{3.0cm}|L{5.1cm}|L{5.6cm}|}
\hline
\rowcolor{gray!15}\textbf{\#} & \textbf{Bước} & \textbf{Mô tả} & \textbf{Kết quả mong đợi} \\ \hline
\endhead
1 & Mở trang Add a new course & Vào Site administration \textgreater{} Courses \textgreater{} Add a new course (hoặc mở \textless{}BASE\_\allowbreak{}URL\textgreater{}/\allowbreak{}course/\allowbreak{}edit.php?category=0). & Form "Add a new course" hiển thị: Course full name, Course short name, Course category (No selection), Course visibility, Course start date, Course end date, Course ID number. \\ \hline
2 & Nhập Course full name & Nhập "TC002 Course 010" vào ô Course full name. & Ô Course full name hiển thị giá trị đã nhập. \\ \hline
3 & Nhập Course short name & Nhập "TC002-010" vào ô Course short name. & Ô Course short name hiển thị giá trị đã nhập. \\ \hline
4 & Chọn Course category & Click ô Search của Course category, gõ "Category 1" và chọn "Category 1" trong danh sách. & Course category hiển thị "Category 1". \\ \hline
5 & Chọn Course visibility & Chọn "Show" trong dropdown Course visibility. & Dropdown hiển thị "Show". \\ \hline
6 & Đặt Course start date & Chọn ngày/\allowbreak{}tháng/\allowbreak{}năm/\allowbreak{}giờ/\allowbreak{}phút = 1 January 2026 08:00. & Course start date hiển thị 1 January 2026 08:00. \\ \hline
7 & Đặt Course end date & Tick "Enable" của Course end date, chọn 1 January 2026 08:00. & Course end date hiển thị 1 January 2026 08:00. \\ \hline
8 & Nhập Course ID number & Để trống (xóa hết nội dung) vào ô Course ID number. & Ô Course ID number hiển thị giá trị đã nhập. \\ \hline
9 & Lưu course & Click nút "Save and display". & Form không được lưu, vẫn ở trang Add a new course; hiển thị lỗi "The course end date must be after the start date.". Không có course mới được tạo. \\ \hline
\end{longtable}
}

\bigskip\noindent\textbf{TC-002-031 -- DT-R4: Thiếu toàn bộ trường bắt buộc}\\
\noindent\textit{Precondition:} Moodle demo (Moodle 5.1, Docker) đang chạy tại http:/\allowbreak{}/\allowbreak{}127.0.0.1:9501 giao diện Moodle và trình duyệt là English; đã đăng nhập bằng admin (admin/\allowbreak{}admin). Đã có category "Category 1". Script Setup đã tạo sẵn course Full name "Existing Course", Short name "EXIST01", ID number "ID-EXIST" trong "Category 1". Chưa tồn tại course nào có Short name/\allowbreak{}ID number dùng trong test case này (script Cleanup xóa các course TC002-* sau khi chạy).
\smallskip

{\footnotesize
\begin{longtable}{|c|L{3.0cm}|L{5.1cm}|L{5.6cm}|}
\hline
\rowcolor{gray!15}\textbf{\#} & \textbf{Bước} & \textbf{Mô tả} & \textbf{Kết quả mong đợi} \\ \hline
\endhead
1 & Mở trang Add a new course & Vào Site administration \textgreater{} Courses \textgreater{} Add a new course (hoặc mở \textless{}BASE\_\allowbreak{}URL\textgreater{}/\allowbreak{}course/\allowbreak{}edit.php?category=0). & Form "Add a new course" hiển thị: Course full name, Course short name, Course category (No selection), Course visibility, Course start date, Course end date, Course ID number. \\ \hline
2 & Nhập Course full name & Để trống (xóa hết nội dung) vào ô Course full name. & Ô Course full name hiển thị giá trị đã nhập. \\ \hline
3 & Nhập Course short name & Để trống (xóa hết nội dung) vào ô Course short name. & Ô Course short name hiển thị giá trị đã nhập. \\ \hline
4 & Không chọn Course category & Giữ nguyên Course category là "No selection". & Course category vẫn là "No selection". \\ \hline
5 & Chọn Course visibility & Chọn "Show" trong dropdown Course visibility. & Dropdown hiển thị "Show". \\ \hline
6 & Đặt Course start date & Chọn ngày/\allowbreak{}tháng/\allowbreak{}năm/\allowbreak{}giờ/\allowbreak{}phút = 1 January 2026 08:00. & Course start date hiển thị 1 January 2026 08:00. \\ \hline
7 & Đặt Course end date & Tick "Enable" của Course end date, chọn 1 January 2027 08:00. & Course end date hiển thị 1 January 2027 08:00. \\ \hline
8 & Nhập Course ID number & Để trống (xóa hết nội dung) vào ô Course ID number. & Ô Course ID number hiển thị giá trị đã nhập. \\ \hline
9 & Lưu course & Click nút "Save and display". & Form không được lưu, vẫn ở trang Add a new course; hiển thị lỗi "Missing full name", "Missing short name" và "You must supply a value here." cùng lúc. Không có course mới được tạo. \\ \hline
\end{longtable}
}

\bigskip\noindent\textbf{TC-002-037 -- UC-P5: Short name trùng rồi đổi}\\
\noindent\textit{Precondition:} Moodle demo (Moodle 5.1, Docker) đang chạy tại http:/\allowbreak{}/\allowbreak{}127.0.0.1:9501 giao diện Moodle và trình duyệt là English; đã đăng nhập bằng admin (admin/\allowbreak{}admin). Đã có category "Category 1". Script Setup đã tạo sẵn course Full name "Existing Course", Short name "EXIST01", ID number "ID-EXIST" trong "Category 1". Chưa tồn tại course nào có Short name/\allowbreak{}ID number dùng trong test case này (script Cleanup xóa các course TC002-* sau khi chạy).
\smallskip

{\footnotesize
\begin{longtable}{|c|L{3.0cm}|L{5.1cm}|L{5.6cm}|}
\hline
\rowcolor{gray!15}\textbf{\#} & \textbf{Bước} & \textbf{Mô tả} & \textbf{Kết quả mong đợi} \\ \hline
\endhead
1 & Mở trang quản lý course & Vào Site administration \textgreater{} Courses \textgreater{} Manage courses and categories, chọn "Category 1". & Trang Course and category management hiển thị, liệt kê course của Category 1. \\ \hline
2 & Mở form tạo course & Click link "Create new course". & Form "Add a new course" hiển thị, Course category mặc định là "Category 1". \\ \hline
3 & Nhập Course full name & Nhập "TC002 UC Dup" vào ô Course full name. & Ô hiển thị giá trị đã nhập. \\ \hline
4 & Nhập Course short name & Nhập "EXIST01" vào ô Course short name. & Ô hiển thị giá trị đã nhập. \\ \hline
5 & Gửi form & Click "Save and display". & Hiển thị "Short name is already used for another course (Existing Course)". \\ \hline
6 & Đổi Short name & Xóa ô Course short name, nhập "TC002-UC5". & Ô hiển thị "TC002-UC5". \\ \hline
7 & Gửi lại & Click "Save and display". & Trang course "TC002 UC Dup" hiển thị. \\ \hline
\end{longtable}
}


